% Reconstructed from the compiled lecture-note PDF.
\chapter{Optimal control from first principles}
\section{From calculus of variations to control}

In the classical calculus of variations one chooses a curve q(t) to minimize


Z tf L(q(t), q′(t)) dt. 0


Optimal control separates the velocity choice from the state. One prescribes


q′(t) = f(q(t), u(t)),


where u(t) is a control taking values in a fixed set U, and minimizes


Z tf l(q(t), u(t)) dt. 0


The control may be only measurable. Discontinuous controls are not pathologies to be removed; they are often the natural optimizers.


\begin{displaytext}
In the Reinhardt problem, \
q = (g, X), U = UT ,
\end{displaytext}

and the control distributes curvature among the synchronized boundary arcs.


\section{Why a multiplier is needed}

The state equation is a constraint. Introduce a time-dependent row vector or covector p(t) and consider the augmented integral


Z tf λcostl(q, u) + p · (f(q, u) −q′) dt, 0


\begin{displaytext}
where λcost ≤0 is a constant multiplier. Integrating the term −p · q′ by parts and varying q suggests
\end{displaytext}

\begin{displaytext}
p′ = −∂H , \
∂q
\end{displaytext}

\begin{displaytext}
where \
H(q, p, u) = p · f(q, u) + λcostl(q, u)
\end{displaytext}

is the control-dependent Hamiltonian. Varying p recovers


\begin{displaytext}
∂H \
q′ = = f(q, u). \
∂p
\end{displaytext}

The remaining variation in u leads not to an ordinary derivative equation when U has corners, but to a pointwise maximization rule.


\section{The Pontryagin maximum principle}

\begin{insightbox}{Standard theorem used as a black box}
\end{insightbox}

\begin{displaytext}
Pontryagin maximum principle, simplified free-time form. Let q∗(t), u∗(t) be an \
optimal trajectory for a smooth control system with compact control set. Then there exist a \
nonzero multiplier pair \
(λcost, p(t)), λcost ≤0,
\end{displaytext}

such that, almost everywhere,


\begin{displaytext}
∂H \
q′ = , \
∂p \
p′ = −∂H , \
∂q \
H(q, p, u∗) = max H(q, p, u). \
u∈U
\end{displaytext}

For a free-terminal-time problem with no terminal cost,


H(q(t), p(t), u∗(t)) = 0


along the extremal. Endpoint constraints give transversality conditions on p.


A state–costate trajectory satisfying these necessary conditions is a Pontryagin extremal. Not every extremal is a minimizer, just as not every critical point in calculus is a minimum.


\section{Normal and abnormal extremals}

\begin{displaytext}
If λcost ̸= 0, scale the multiplier so that
\end{displaytext}

\begin{displaytext}
λcost = −1.
\end{displaytext}

\begin{displaytext}
The extremal is then normal. If \
λcost = 0,
\end{displaytext}

the cost disappears from the first-order equations; the extremal is abnormal. Abnormal extremals reflect the geometry of the constraints rather than the cost itself. They can be genuine, but they can also be artifacts of a redundant formulation.


The Reinhardt problem has anomalous abnormal edge extremals in its full encoding. The source therefore introduces a reduced edge-control problem to remove them when proving finite switching along an edge.


\section{Cotangent vectors without differential geometry}

\begin{displaytext}
At a point q ∈Rn, a tangent vector is an ordinary vector v ∈Rn. A cotangent vector is a linear \
functional \
p : v 7→p · v.
\end{displaytext}

\begin{displaytext}
For a matrix group, left translation identifies every tangent space with the same matrix space. Thus \
the costates for g ∈SL2(R) and X ∈sl2(R) can both be represented by traceless matrices, using \
the nondegenerate trace pairing \
⟨A, B⟩= tr(AB).
\end{displaytext}

\begin{displaytext}
This is why the source writes the costates as matrices Λ1 and Λ2.
\end{displaytext}

\section{Maximized Hamiltonian}

\begin{displaytext}
Define \
H+(q, p) = max H(q, p, u). \
u∈U
\end{displaytext}

If the maximizer is unique and depends smoothly on (q, p), then the extremal follows the ordinary Hamiltonian flow of H+. When the maximizer jumps from one extreme point of U to another, H+ is not differentiable across the switching surface, but the state and costate remain continuous.


\begin{statementbox}{Definition}
Definition 8.1 (Bang-bang control). A control is bang-bang if it takes values only at extreme points of the control set, except possibly on a set of measure zero.
\end{statementbox}

For an interval control U = [−1, 1], a bang-bang control takes the values −1 and 1. For the Reinhardt triangle, it takes the three vertex values e0, e1, e2.


\section{Switching functions}

Suppose the Hamiltonian evaluated at two candidate controls is Hi(t) and Hj(t). Their difference


\begin{displaytext}
χij(t) = Hi(t) −Hj(t)
\end{displaytext}

\begin{displaytext}
is a switching function. The maximization rule chooses i when χij > 0 and j when χij < 0. A \
switch can occur only at a zero of χij.
\end{displaytext}

\begin{displaytext}
If a zero is simple, \
χij(t0) = 0, χ′ij(t0) ̸= 0,
\end{displaytext}

the switching time varies smoothly with initial conditions by the implicit-function theorem. Multiple zeros are the source of discontinuities in Poincaré maps and require discriminants and resultants later.


\section{Faces of a convex control set}

A subset F of a convex set U is a face if whenever an interior point of a segment in U lies in F, both endpoints lie in F. The faces of a triangle are its vertices, edges, and the whole triangle.


In the Reinhardt Hamiltonian, the control-dependent term is a ratio of two affine functions of u. Along a line segment u(t) = tu1 + (1 −t)u0, such a ratio has the form at + b ct + d. Its derivative has constant sign unless it is identically constant. Therefore the set of maximizers is always a face of UT .


Consequently, at each time exactly one of the following occurs:


1. a unique vertex maximizes: regular bang-bang motion; 2. an entire edge maximizes: an edge singularity; 3. the whole triangle maximizes: the fully singular locus.


\section{Singular arcs}

A singular arc is a time interval on which the maximization rule fails to determine a unique control. The name refers to singularity of the first-order necessary conditions, not necessarily to a geometric singularity of the state curve.


In the Reinhardt problem, full-triangle singularity forces the circular state. Edge singularities are handled by the auxiliary edge-control problem. Outside the full singular locus, edge extremality will imply finite bang-bang switching.


\section{Chattering}

A bang-bang control may have infinitely many switches in a finite interval. If the switching times


t1 < t2 < · · · < tn < · · ·


converge to a finite limit t∗, the control chatters. The state can converge smoothly even though the control oscillates infinitely rapidly.


Chattering cannot be excluded by compactness or first-order optimality alone. Fuller’s classical example shows that it can be the genuine optimal behavior. The central difficulty of the source paper is to prove that a periodic Reinhardt minimizer does not chatter into the circular singular locus.


\section{Poincaré maps}

For a bang-bang flow, record only the state at switching times. If qn is the state at the nth switch, define qn+1 = F(qn).


This F is a Poincaré first-recurrence map. It converts a continuous-time control system into a discrete dynamical system.


Near a chattering singularity, the time intervals shrink, but after a suitable blowup the angular switching data can converge to a nontrivial fixed point of F. This is the key idea used in Part V.


\begin{insightbox}{Chapter summary}
\end{insightbox}

The Pontryagin maximum principle augments the state with a costate and replaces minimization by Hamiltonian evolution plus pointwise control maximization. For the triangular control set, maximizers are faces. Unique vertex maximizers give bang-bang motion; edges and the whole triangle give singular behavior. Infinite vertex switching in finite time is chattering, the only obstruction that survives the first-order analysis.

